% Generated draft: requires mathematical and source review.
\documentclass[11pt]{article}
\usepackage[margin=25mm]{geometry}
\usepackage{fontspec}
\setmainfont{Latin Modern Roman}
\usepackage{amsmath,amssymb,mathtools,unicode-math}
\setmathfont{Latin Modern Math}
\usepackage{enumitem,needspace,xurl,hyperref}
\setlength{\parindent}{0pt}
\setlength{\parskip}{5pt}
\setlength{\emergencystretch}{3em}
\newcommand{\N}{\mathbb{N}}
\newcommand{\Z}{\mathbb{Z}}
\newcommand{\Q}{\mathbb{Q}}
\newcommand{\R}{\mathbb{R}}
\newcommand{\C}{\mathbb{C}}
\newcommand{\F}{\mathbb{F}}
\newcommand{\A}{\mathbb{A}}
\newcommand{\PP}{\mathbb{P}}
\newcommand{\E}{\mathbb{E}}
\newcommand{\eps}{\varepsilon}
\newcommand{\dd}{\,\mathrm{d}}
\begin{document}
\section*{UnsolvedMath ID: 30005782}
\textbf{Source problem code:} catalogue entry
\subsection{Definitions and mathematical statement}
Let \(S\) be a set and let \(M=(m_{st})_{s,t\in S}\) be a symmetric matrix with \(m_{ss}=1\) and \(m_{st}\in\{2,3,4,\ldots\}\cup\{\infty\}\) for distinct \(s,t\). For a positive integer \(m\), let \(w_m(s,t)\) denote the alternating word of length \(m\) beginning with \(s\). The associated Artin group is
\[
A(M)=\left\langle S\;\middle|\;w_{m_{st}}(s,t)=w_{m_{st}}(t,s)\quad(s\ne t,\ m_{st}<\infty)\right\rangle.
\]
An entry \(m_{st}=\infty\) imposes no relation between the corresponding generators. The group has finite rank when \(S\) is finite.

The recoverable torsion-freeness question in the supplied introductory passage is whether every Artin group is torsion-free: for every matrix \(M\) in the intended class,
\[
\forall g\in A(M)\ \forall n\in\Z_{\geq 1},\qquad g^n=1\ \Longrightarrow\ g=1.
\]
The supplied passage does not explicitly specify whether its class of Artin groups is restricted to finite rank; the definition above allows either convention, and this scope qualification remains unresolved.

The supplied dated literature assessment additionally identifies solvability of the word problem as an intended question. For a specified finite set \(S\) and matrix \(M\), the conventional formulation is: does there exist an algorithm which, on every input finite word \(u\) in the generators \(S\) and their formal inverses, halts and correctly decides whether \(u=1\) in \(A(M)\)? The corresponding assertion for all finite-rank Artin groups has quantifier order
\[
\forall\text{ finite-rank Artin presentations }M\quad
\exists\text{ an algorithm deciding the word problem in }A(M).
\]
This does not by itself assert the existence of a single algorithm taking both \(M\) and \(u\) as input. The supplied record does not specify uniformity, nor does it specify an effective presentation convention for an infinite generating set.

These are qualified formulations of the questions mentioned in the record, not a recovered complete standalone problem. The original passage terminates at the words ``whether they a have solvable'', and its continuation is missing. The word-problem interpretation comes from the supplied literature assessment. The record labels the statement \texttt{unrecoverable} and describes it as introductory workshop context duplicating concrete Artin-group problems elsewhere; those further problems are not supplied here.
\subsection{Short English statement}
The introduction asks whether Artin groups have no nonidentity elements of finite order. The supplied curation also identifies the question of whether equality to the identity can always be decided algorithmically from a word. The passage is truncated and marked \texttt{unrecoverable}, so its precise scope and complete intended target require review. Its supplied open-status assessment is dated 2026-08-22 and has not been independently verified.
\subsection{Sources}
\begin{itemize}
\item Local UnsolvedMath registry, \texttt{data/unsolvedmath/problems.json}, supplied record with permanent ID \(30005782\), exact source code \texttt{catalogue entry}, title ``Artin-Group Workshop Introduction'', in the Oberwolfach Reports Open Problems collection. The supplied record gives proposed year 2024 and no individual proposer.
\item Supplied report citation: \emph{Mini-Worskhop: Artin Groups meet Triangulated Categories} (2024), Oberwolfach Reports, DOI 10.4171/owr/2024/4, pp.\ 203--234. Source URL and literature-source link labelled ``original OWR report'': \url{https://doi.org/10.4171/owr/2024/4}. This is the authoritative report citation supplied by the record; it has not been independently consulted for this draft.
\item Supplied literature assessment, checked at exactly \texttt{2026-08-22}: torsion-freeness and solvability of the word problem are described as established for substantial families of Artin groups but open for general Artin groups as stated in the report. This is supplied dated curation, not an independently verified assessment of the literature. The background expressly requests independent verification of time-sensitive status and cited claims.
\item Supplied statement-verification note: the official report is said to confirm that this passage is introductory workshop context duplicating concrete Artin-group problems preserved elsewhere. This is a supplied verification claim, not a source visit performed for this draft. The record separately assigns statement status \texttt{unrecoverable} and problem status \texttt{open}.
\item Contextual references retained from the passage: [Dav83], associated there with Davis' construction of a counterexample to the claim that all aspherical manifolds have Euclidean universal covers; and [Ago13], associated there with Agol's proof of the virtual Haken conjecture. Complete bibliographic entries for these labels are not supplied, and these contextual claims have not been independently verified.
\end{itemize}
\end{document}
